% =============================================================================
%  Notas del Profesor — Sesión 05: Bit Manipulation
%  Programación Avanzada — Universidad Panamericana
%  Compilar: pdflatex sesion05_bit_manipulation_profesor.tex
% =============================================================================
\documentclass[12pt, oneside]{article}
\usepackage[profesor]{progavanzada}

\begin{document}

\portada{Sesión 05}{Bit Manipulation}{2026}

\tableofcontents
\newpage

% =============================================================================
\section{Objetivos de la sesión}
% =============================================================================

Al finalizar la sesión el alumno será capaz de:

\begin{enumerate}
  \item Leer y escribir literales binarios (\cpp{0b}) y hexadecimales (\cpp{0x})
        en C++.
  \item Usar \texttt{bitset<N>} para visualizar la representación binaria de
        un valor.
  \item Calcular el resultado de \cpp{<<}, \cpp{>>},
        \texttt{\textasciitilde}, \texttt{\&}, \texttt{|} y
        \texttt{\textasciicircum} sobre valores de 8 bits a mano.
  \item Identificar los tres patrones especiales (todo ceros, todo unos, máscara
        de un bit) y reconocerlos en código.
  \item Evaluar expresiones que combinen varios operadores de bits paso a paso.
\end{enumerate}

\begin{notaprofesor}[Duración estimada]
  50 minutos. Sugerido: 5 min apertura + 10 min visualización y literales +
  10 min patrones especiales + 15 min operadores lógicos + 10 min ejercicios.
  La sesión 06 es práctica en el cuaderno; estos ejercicios son en papel.
\end{notaprofesor}

% =============================================================================
\section{Secuencia de clase}
% =============================================================================

\subsection{Apertura (5 min)}

Arranca con esta pregunta al grupo:

\begin{notaprofesor}[Pregunta de arranque]
  \textit{``¿Cuántas operaciones necesita el procesador para calcular
  \texttt{x * 8}?''}\\[4pt]
  La mayoría dirá ``una multiplicación''. La respuesta real: si usas
  \texttt{x << 3}, el procesador lo resuelve con \textbf{un solo ciclo de
  reloj}. Una multiplicación normal tarda entre 3 y 5 ciclos.
  Ese ahorro, repetido millones de veces, importa.\\[4pt]
  Sigue con: \textit{``Hoy vamos a aprender a hablarle al procesador en su
  idioma natural: bits.''}
\end{notaprofesor}

\subsection{Visualización: \texttt{bitset} y literales (10 min)}

Muestra en el cuaderno (o proyecta) la celda de código con \texttt{bitset}:
\begin{verbatim}
int x = 30;
cout << bin8(x);   // 00011110
\end{verbatim}

Pregunta al grupo: \textit{``¿Por qué \texttt{00011110}?''} Deja que los
alumnos cuenten las potencias de 2 en voz alta: $16+8+4+2=30$.

\begin{notaprofesor}[Literal binario vs.\ decimal]
  Muestra \texttt{0b00011110}, \texttt{0x1E} y \texttt{30} imprimiendo
  exactamente lo mismo. El punto: la base es solo presentación. El compilador
  guarda los mismos bits sin importar cómo escribiste el literal.\\[4pt]
  Pregunta de sondeo: \textit{``¿Qué imprimiría
  \texttt{cout << 0b0101}?''} — Respuesta: 5. No imprime ``0101''; para eso
  se necesita \texttt{bitset}.
\end{notaprofesor}

Después introduce los patrones especiales. La tabla de \cpp{1 << k}
funciona bien en el cuaderno porque los alumnos pueden ejecutarla celda
por celda y ver el bit moverse.

\subsection{Operadores bit a bit (15 min)}

Presenta NOT, AND, OR y XOR con sus tablas de verdad. Para cada uno:
\begin{enumerate}
  \item Tabla de verdad de 1 bit.
  \item Ejemplo con 8 bits (alinear operandos en columnas en el pizarrón).
  \item El ``uso típico'' — para qué sirve en la práctica.
\end{enumerate}

\begin{notaprofesor}[Confusión frecuente: \texttt{\&} vs.\ \texttt{\&\&}]
  Varios alumnos confundirán \texttt{\&} (AND bit a bit) con \texttt{\&\&}
  (AND lógico). El AND lógico trata todo valor distinto de cero como
  \texttt{true}; el AND bit a bit opera posición a posición. Ejemplo
  que los distingue:
  \begin{center}
  \begin{tabular}{lccc}
    & \texttt{a = 2} & \texttt{b = 1} & Resultado \\
    \hline
    \texttt{a \&\& b} & (true) & (true) & \texttt{1} (true) \\
    \texttt{a \& b}   & \texttt{10} & \texttt{01} & \texttt{0} \\
  \end{tabular}
  \end{center}
  El mismo par de valores da resultados completamente distintos.
\end{notaprofesor}

\begin{notaprofesor}[El XOR como cifrador]
  Si el grupo lleva bien el ritmo, menciona el truco: aplicar XOR con
  la misma clave dos veces recupera el original (\texttt{(x \^{} k) \^{} k = x}).
  Es la base del cifrado Vernam (one-time pad), el único sistema de cifrado
  matemáticamente irrompible. Funciona bien como anécdota de cierre.
\end{notaprofesor}

\subsection{Ejercicios en clase (10 min)}

Los ejercicios de la sección 7 de las notas del alumno. Recomendado:
\begin{itemize}
  \item Los ejercicios 1, 4 y 7a en clase (desplazamientos + tabla AND/OR
        + expresión combinada).
  \item Los demás quedan como tarea para reforzar antes de la sesión 06
        (práctica en cuaderno).
\end{itemize}

\begin{notaprofesor}[Ejercicio 9: \texttt{n \& (n-1)}]
  Este ejercicio es una pista hacia la sesión 06. \texttt{n \& (n-1)}
  apaga el bit encendido de menor peso. Si \texttt{n} es potencia de 2,
  solo tiene un bit encendido, así que el resultado es cero.
  Es el corazón del algoritmo para contar bits activos (popcount eficiente)
  que verán en la sesión 06.
\end{notaprofesor}

% =============================================================================
\section{FAQs}
% =============================================================================

\begin{notaprofesor}[¿Por qué \texttt{\textasciitilde 30 = -31}?]
  En complemento a 2, \texttt{\textasciitilde x = -(x+1)}.
  Intuitivamente: \texttt{x + \textasciitilde x = -1} (todos los bits
  encendidos), así que \texttt{\textasciitilde x = -1 - x = -(x+1)}.\\[4pt]
  Verificación rápida: \texttt{\textasciitilde 0 = -1} (todo 1s),
  \texttt{\textasciitilde(-1) = 0} (todo 0s).
\end{notaprofesor}

\begin{notaprofesor}[¿Los desplazamientos son más rápidos que multiplicar?]
  En hardware moderno la diferencia es menor que antes (los procesadores
  tienen multiplicadores muy optimizados), pero el compilador \textbf{ya}
  convierte \texttt{x * 8} en \texttt{x << 3} automáticamente cuando puede.
  La razón real para aprender desplazamientos: leer código de sistemas y
  programación competitiva donde el idioma son bits, no aritmética.
\end{notaprofesor}

\begin{notaprofesor}[¿Qué pasa con \texttt{>>} en números negativos?]
  El estándar C++ (hasta C++19) dice que el comportamiento de \texttt{>>}
  sobre negativos es \textit{implementation-defined}. En la práctica, g++
  y clang hacen extensión de signo (aritmética), así que \texttt{-8 >> 1 = -4}.
  Pero el código correcto usa \texttt{unsigned int} si necesita desplazar
  a la derecha de forma portable.
\end{notaprofesor}

\begin{notaprofesor}[¿Por qué \texttt{-1} tiene todos los bits en 1?]
  En complemento a 2, $-1$ se almacena como $2^n - 1$. Para 32 bits:
  $2^{32} - 1 = 4{,}294{,}967{,}295 = \texttt{0xFFFFFFFF}$ = todo 1s.\\[4pt]
  Forma alternativa de verlo: \texttt{0} tiene todo en 0. Le restamos 1
  y como no podemos (prestamos de un bit que no existe), todos los bits
  se encadenan y quedan en 1.
\end{notaprofesor}

% =============================================================================
\section{Soluciones a los ejercicios}
% =============================================================================

\subsection{Desplazamientos}

\begin{notaprofesor}[Ejercicio 1]
\begin{enumerate}[label=\alph*)]
  \item \texttt{7 << 2} $= 7 \times 4 = 28$; binario: \texttt{00011100}.
  \item \texttt{100 >> 2} $= \lfloor 100/4 \rfloor = 25$; binario: \texttt{00011001}.
  \item \texttt{1 << 5} $= 32$; binario: \texttt{00100000}.
  \item \texttt{96 >> 3} $= \lfloor 96/8 \rfloor = 12$; binario: \texttt{00001100}.
\end{enumerate}
\end{notaprofesor}

\begin{notaprofesor}[Ejercicio 2]
  $48 / 3 = 16 = 2^4$, así que cuatro desplazamientos a la izquierda.
  Verificación: \texttt{3 << 4} $= 3 \times 16 = 48$. \checkmark
\end{notaprofesor}

\begin{notaprofesor}[Ejercicio 3]
  Si \texttt{x} tiene a lo más 8 bits, desplazarlo 8 posiciones expulsa
  todos sus bits; el resultado es \texttt{0}. (Técnicamente, si \texttt{x}
  es \texttt{int} de 32 bits, los 24 bits altos son 0, y todos los bits
  de \texttt{x} ``caen'' fuera por la derecha.)
\end{notaprofesor}

\subsection{NOT, AND, OR y XOR}

\begin{notaprofesor}[Ejercicio 4]
\begin{enumerate}[label=\alph*)]
  \item \texttt{17 \& 12}:\\
        \texttt{00010001 \& 00001100 = 00000000 = 0}.
        No comparten ningún bit encendido.
  \item \texttt{13 | 12}:\\
        \texttt{00001101 | 00001100 = 00001101 = 13}.
        Los bits de 12 ya están en 13, así que OR no aporta nada nuevo.
  \item \texttt{17 \^{} 12}:\\
        \texttt{00010001 \^{} 00001100 = 00011101 = 29}.
  \item \texttt{\textasciitilde 12 = -(12+1) = -13}.
        En 32 bits: \texttt{11111111111111111111111111110011}.
\end{enumerate}
\end{notaprofesor}

\begin{notaprofesor}[Ejercicio 5 — \texttt{x = 0b10110011}]
\begin{enumerate}[label=\alph*)]
  \item \texttt{x \& 0b00001111 = 0b00000011 = 3}.
        AND con \texttt{0x0F} \textbf{conserva solo los 4 bits bajos}
        (nibble bajo) y apaga los 4 altos. Operación clásica para extraer
        el nibble bajo.
  \item \texttt{x | 0b00001111 = 0b10111111 = 191}.
        OR con \texttt{0x0F} \textbf{enciende los 4 bits bajos} sin tocar
        los demás.
  \item \texttt{x \^{} 0b11111111 = 0b01001100 = 76}.
        XOR con \texttt{0xFF} \textbf{invierte todos los bits}, equivalente
        a NOT para los 8 bits bajos. Notar que difiere de \texttt{\textasciitilde x}
        en que este solo actúa sobre 8 bits, no los 32.
\end{enumerate}
\end{notaprofesor}

\begin{notaprofesor}[Ejercicio 6 — tabla AND/OR]
\begin{center}
\begin{tabular}{cccc}
  \toprule
  \texttt{a} & \texttt{b} & \texttt{a \& b} & \texttt{a | b} \\
  \midrule
  \texttt{11001100} & \texttt{10101010} & \texttt{10001000} & \texttt{11101110} \\
  \texttt{11110000} & \texttt{00001111} & \texttt{00000000} & \texttt{11111111} \\
  \texttt{10101010} & \texttt{10101010} & \texttt{10101010} & \texttt{10101010} \\
  \bottomrule
\end{tabular}
\end{center}
La segunda fila es interesante: cuando los patrones son \textbf{complementos}
entre sí, AND da todo 0s y OR da todo 1s.
La tercera fila: AND y OR con sí mismo devuelve el original.
\end{notaprofesor}

\subsection{Combinando operadores}

\begin{notaprofesor}[Ejercicio 7a — \texttt{(\textasciitilde 7 | (1<<1) | (1<<3) | (1<<5)) \& 29}]
  Paso a paso:
  \begin{enumerate}
    \item \texttt{\textasciitilde 7 = ...11111000} (apaga los bits 0, 1, 2)
    \item \texttt{| (1<<1) = ...11111010} (enciende el bit 1)
    \item \texttt{| (1<<3) = ...11111010} (bit 3 ya estaba)
    \item \texttt{| (1<<5) = ...11111010} (bit 5 ya estaba)
    \item \texttt{29 = 0b00011101}
    \item AND: \texttt{0b00011101 \& 0b11111010 = 0b00011000 = 24}
  \end{enumerate}
\end{notaprofesor}

\begin{notaprofesor}[Ejercicio 7b — \texttt{(\textasciitilde(13 >> 2)) \& 39}]
  \begin{enumerate}
    \item \texttt{13 = 0b00001101}
    \item \texttt{13 >> 2 = 0b00000011 = 3}
    \item \texttt{\textasciitilde 3 = ...11111100}
    \item \texttt{39 = 0b00100111}
    \item \texttt{\textasciitilde 3 \& 39 = 0b00100100 = 36}
  \end{enumerate}
\end{notaprofesor}

\begin{notaprofesor}[Ejercicio 8 — \texttt{x \^{} y \^{} y}]
  Por las propiedades del XOR: \texttt{y \^{} y = 0}, y \texttt{x \^{} 0 = x}.
  Resultado: siempre \texttt{x}, sin importar \texttt{y}.
  Este truco es la base del \textit{XOR swap} (intercambiar dos variables
  sin una temporal) que verán en la sesión 06.
\end{notaprofesor}

\begin{notaprofesor}[Ejercicio 9 — \texttt{n \& (n-1)} con \texttt{n = 12}]
  \texttt{12 = 0b00001100}. \quad \texttt{11 = 0b00001011}.\\[4pt]
  \texttt{12 \& 11 = 0b00001000 = 8}.\\[4pt]
  El resultado apaga el bit encendido de \textbf{menor peso} de \texttt{n}
  (el bit 2 en este caso). Si \texttt{n} es potencia de 2 (un solo bit
  encendido), el resultado es 0. Esta es la prueba más eficiente para
  saber si un número es potencia de 2: \texttt{n > 0 \&\& (n \& (n-1)) == 0}.
\end{notaprofesor}

% =============================================================================
\section{Demo en vivo sugerida}
% =============================================================================

\begin{notaprofesor}[Conversión mayúsculas/minúsculas con XOR]
  Esta demo conecta la sesión 04 (ASCII) con la sesión 05 (bits):
  \begin{verbatim}
  char c = 'A';              // 0x41 = 0100 0001
  char lower = c ^ 0x20;    // 0x61 = 0110 0001  = 'a'
  char back  = lower ^ 0x20; // 0x41 = 0100 0001  = 'A'
  \end{verbatim}
  El XOR con \texttt{0x20} invierte el bit 5, que es exactamente el bit
  que distingue mayúsculas de minúsculas en ASCII.
  Aplicarlo dos veces devuelve la letra original.\\[4pt]
  Punto de discusión: ¿funciona igual con \texttt{|} o \texttt{\&}?
  No: OR solo sube, AND solo baja. XOR es el único que \textit{alterna}.
\end{notaprofesor}

\begin{notaprofesor}[Sugerencia de apertura del cuaderno]
  Si la sesión 05 va bien y sobra tiempo, abre el cuaderno y ejecuta la
  celda de los patrones especiales (\texttt{1 << k} del 0 al 7). Ver el bit
  ``caminar'' hacia la izquierda en la terminal suele generar buena discusión
  sobre lo que físicamente representa cada bit en la memoria.
\end{notaprofesor}

\end{document}
